% GENERATED FILE. Edit canonical lessons, then run npm run generate:books.
% canonical-source-hash: fnv1a64:87ad329b3e59b93d
% canonical-lessons: PT-C04-adeus, PT-C04-ate-logo, PT-C04-ate-amanha, PT-C04-ate-breve, PT-C04-practice

\chapter{Farewells}
\label{ch:pt-farewells}
\hlchaptermodality{4}

\begin{chapteropening}
\textbf{By the end of this chapter:} \emph{I can close a Portuguese conversation and pick the goodbye that matches when I will next see the person.}

Open with Olá, ask Tudo bem?, then choose the exit — até logo, até amanhã, até breve, or adeus, to God, for the one that means for good.
\end{chapteropening}

\section[adeus]{adeus — \textquotedblleft{}goodbye,\textquotedblright{} i.e. \textquotedblleft{}to God\textquotedblright{}}

\label{lesson:PT-C04-adeus}



\begin{quote}
You can greet, ask \emph{tudo bem?}, and reply. Now close the conversation. The Portuguese goodbye is \textbf{adeus} — and it's a tiny prayer: \textbf{\textquotedblleft{}to God.\textquotedblright{}}
\end{quote}

\begin{sounds}
\begin{itemize}
\raggedright
  \item \texttt{nasal-eu} — \textbf{adeus} $\approx$ \emph{ah-\textbf{DEH}-oosh}: the \emph{eu} is a gliding vowel; final \emph{-s} is a soft \emph{sh} in Portugal (\emph{-s} $\to$ \emph{ss} in Brazil). Stress the \emph{-deus}.
\end{itemize}
\end{sounds}

\begin{cousinweb}
\textbf{adeus} = \textbf{a} (\textquotedblleft{}to\textquotedblright{}) + \textbf{Deus} (\textquotedblleft{}God\textquotedblright{}), from Latin \textbf{ad Deum}, \emph{\textquotedblleft{}to God.\textquotedblright{}} It's the fossil of a blessing — \emph{\textquotedblleft{}{[}I commend you{]} to God\textquotedblright{}} — said at parting, exactly like the English \textbf{goodbye}, which is a worn-down \emph{\textquotedblleft{}God be with ye.\textquotedblright{}}

This is one of the most beautiful family resemblances anywhere — five languages, one farewell prayer:

\noindent
\begin{tabularx}{\linewidth}{@{}>{\raggedright\arraybackslash}X>{\raggedright\arraybackslash}X>{\raggedright\arraybackslash}X@{}}
\toprule
\textbf{language} & \textbf{goodbye} & \textbf{=} \\
\midrule
\textbf{Portuguese} & \emph{adeus} & a + Deus \\
\textbf{Spanish} & \emph{adiós} & a + Dios \\
\textbf{French} & \emph{adieu} & à + Dieu \\
\textbf{Italian} & \emph{addio} & a + Dio \\
\textbf{English} & \emph{goodbye} & \textquotedblleft{}God be with ye\textquotedblright{} \\
\bottomrule
\end{tabularx}

The \emph{Deus} here is the same one in English \textbf{deity}, \textbf{divine}, \textbf{adieu}, and Spanish \emph{Dios}. When you say \emph{adeus}, you're handing someone to God for safekeeping until next time.
\end{cousinweb}

\subsection*{Your turn}
Say these aloud:

\begin{itemize}
\raggedright
  \item \textquotedblleft{}adeus\textquotedblright{} — \emph{ah-DEH-oosh}
  \item adeus / adiós / adieu / addio — one \textquotedblleft{}to God\textquotedblright{} across four languages
  \item even English \textquotedblleft{}goodbye\textquotedblright{} = \textquotedblleft{}God be with ye\textquotedblright{}
\end{itemize}

\begin{tcolorbox}[breakable,colback=teal!4,colframe=teal!35!black,title={Before you move on}]
What does \emph{adeus} literally mean, and from what? (\textquotedblleft{}To God,\textquotedblright{} $\leftarrow$ \emph{ad Deum}.) Name two sister-language goodbyes with the same \textquotedblleft{}to God\textquotedblright{} origin. (\emph{adiós}, \emph{adieu}, \emph{addio}.) What English goodbye hides \textquotedblleft{}God be with ye\textquotedblright{}? (\emph{Goodbye}.) Next: the softer \textquotedblleft{}see you\textquotedblright{} goodbyes — starting with \textbf{até logo}.
\end{tcolorbox}
\section[até logo]{até logo — \textquotedblleft{}see you later,\textquotedblright{} and a word straight from Arabic}

\label{lesson:PT-C04-ate-logo}



\begin{quote}
The everyday Portuguese goodbye isn't \emph{adeus} (that's a bit final) — it's \textbf{até logo}, \textquotedblleft{}until soon.\textquotedblright{} And its first word carries the same surprise Spanish's \emph{hasta} does: it came from \textbf{Arabic}.
\end{quote}

\begin{sounds}
\begin{itemize}
\raggedright
  \item \textbf{até} = \emph{ah-\textbf{TEH}}, stress the end (the accent on \emph{é}).
  \item \texttt{final-o-is-u} — \textbf{logo} = \emph{LOH-goo}, final \emph{-o} $\to$ \emph{u}.
\end{itemize}
\end{sounds}

\begin{cousinweb}
\begin{itemize}
\raggedright
  \item \textbf{até} = \textquotedblleft{}until, up to.\textquotedblright{} Like Spanish \textbf{hasta}, it comes from \textbf{Arabic \emph{ḥattā} (\ptar{حتى})}, \textquotedblleft{}until\textquotedblright{} — carried into Iberian speech across the \textbf{800 years of al-Andalus}. Both Iberian languages borrowed the \emph{same} Arabic word for this everyday preposition. That's remarkable: languages freely borrow \emph{nouns} (Portuguese has thousands from Arabic — \emph{almofada}, \emph{açúcar}, \emph{azeite}), but a \textbf{function word} like \textquotedblleft{}until\textquotedblright{} almost never crosses. \emph{até} and \emph{hasta} are twin survivors of that deep Arabic layer under Iberian Romance.
  \item \textbf{logo} = \textquotedblleft{}soon, right away, then,\textquotedblright{} from Latin \textbf{locō}, \emph{\textquotedblleft{}in {[}that{]} place\textquotedblright{}} $\to$ \textquotedblleft{}on the spot\textquotedblright{} $\to$ \textquotedblleft{}soon.\textquotedblright{} (Spanish did the identical thing: \textbf{luego} $\leftarrow$ \emph{locō}, and \emph{hasta luego} is the word-for-word twin of \emph{até logo}.)
\end{itemize}

So \textbf{até logo} = \textquotedblleft{}\textbf{until soon}\textquotedblright{} — the exact mirror of Spanish \textbf{hasta luego}, Arabic preposition and all.
\end{cousinweb}

\subsection*{Your turn}
Say these aloud:

\begin{itemize}
\raggedright
  \item \textquotedblleft{}até logo\textquotedblright{} — \emph{ah-TEH LOH-goo}
  \item Portuguese \emph{até} and Spanish \emph{hasta} — both from Arabic \emph{ḥattā}
  \item \emph{até logo} = \emph{hasta luego}, twin for twin
\end{itemize}

\begin{tcolorbox}[breakable,colback=teal!4,colframe=teal!35!black,title={Before you move on}]
Where does \emph{até} come from, and which Spanish word is its twin? (Arabic \emph{ḥattā} — Spanish \emph{hasta}.) Why is borrowing it surprising? (It's a \emph{function word}; languages rarely borrow those.) What does \emph{logo} share with Spanish \emph{luego}? (Both $\leftarrow$ Latin \emph{locō}, \textquotedblleft{}in place\textquotedblright{} $\to$ \textquotedblleft{}soon\textquotedblright{}.) Next: \textbf{até amanhã}.
\end{tcolorbox}
\section[até amanhã]{até amanhã — \textquotedblleft{}see you tomorrow\textquotedblright{}}

\label{lesson:PT-C04-ate-amanha}



\begin{quote}
Same frame — \emph{até} + a “when.” \textbf{Até amanhã} means “until tomorrow.” Its second word is the Portuguese cousin of a Spanish word you already know, wearing a heavy nasal ending.
\end{quote}

\begin{sounds}
\begin{itemize}
\raggedright
  \item \texttt{nasal-anha} — \textbf{amanhã} = \emph{ah-mah-\textbf{NYÃ}}: the \textbf{nh} is the \emph{ny} of \textquotedblleft{}canyon\textquotedblright{} (Portuguese spells the \emph{ñ} sound as \emph{nh}), and the final \textbf{-ã} is a strong \textbf{nasal} \emph{a} — air through the nose, no hard consonant. Stress the end.
\end{itemize}
\end{sounds}

\begin{cousinweb}
\begin{itemize}
\raggedright
  \item \textbf{até} = \textquotedblleft{}until\textquotedblright{} ($\leftarrow$ Arabic \emph{ḥattā}, from the last lesson).
  \item \textbf{amanhã} = \textquotedblleft{}tomorrow,\textquotedblright{} from Latin \textbf{ad māneāna}, \emph{\textquotedblleft{}to the morning {[}hour{]}.\textquotedblright{}}
\end{itemize}

That \textbf{māne} (\textquotedblleft{}morning\textquotedblright{}) root is shared right across Romance — you've now met three of its children:

\begin{itemize}
\raggedright
  \item Portuguese \textbf{amanhã} (\textquotedblleft{}tomorrow\textquotedblright{}) $\leftarrow$ \emph{ad māneāna}.
  \item Spanish \textbf{mañana} — morning \emph{and} tomorrow (the \emph{ñ} spelled with the tilde).
  \item Italian \textbf{domani} (\textquotedblleft{}tomorrow\textquotedblright{}) $\leftarrow$ \emph{dē māne}.
\end{itemize}

Same idea everywhere: tomorrow is \textbf{\textquotedblleft{}the morning to come.\textquotedblright{}} Portuguese writes the \emph{ny} sound as \textbf{nh} (\emph{amanhã}), where Spanish uses \textbf{ñ} (\emph{mañana}) — two spellings, one sound, both from that doubled-\emph{n}.
\end{cousinweb}

\subsection*{Your turn}
Say these aloud:

\begin{itemize}
\raggedright
  \item \textquotedblleft{}até amanhã\textquotedblright{} — \emph{ah-TEH ah-mah-NYÃ}, nasal ending
  \item amanhã / mañana / domani — three children of \emph{māne}, \textquotedblleft{}morning\textquotedblright{}
  \item Portuguese \emph{nh} = Spanish \emph{ñ} = the \emph{ny} sound
\end{itemize}

\begin{tcolorbox}[breakable,colback=teal!4,colframe=teal!35!black,title={Before you move on}]
What does \emph{amanhã} come from, and literally mean? (\emph{Ad māneāna} — \textquotedblleft{}to the morning.\textquotedblright{}) How does Portuguese spell the \emph{ny} sound that Spanish writes as \emph{ñ}? (As \emph{nh}.) Next: \textquotedblleft{}see you soon\textquotedblright{} — \textbf{até breve}.
\end{tcolorbox}
\section[até breve]{até breve — \textquotedblleft{}see you soon\textquotedblright{}}

\label{lesson:PT-C04-ate-breve}



\begin{quote}
The last of the \emph{até} goodbyes: \textbf{até breve}, \textquotedblleft{}until soon\textquotedblright{} — literally \textquotedblleft{}until {[}a{]} \emph{short} {[}while{]}.\textquotedblright{} \emph{Breve} is a word English wears in several coats.
\end{quote}

\begin{sounds}
\begin{itemize}
\raggedright
  \item \textbf{até breve} = \emph{ah-TEH \textbf{BREH}-vee}: open \emph{e} in \emph{breve}, final \emph{-e} softens toward \emph{ee}.
\end{itemize}
\end{sounds}

\begin{cousinweb}
\begin{itemize}
\raggedright
  \item \textbf{até} = \textquotedblleft{}until\textquotedblright{} ($\leftarrow$ Arabic \emph{ḥattā}).
  \item \textbf{breve} = \textquotedblleft{}short, brief,\textquotedblright{} from Latin \textbf{brevis}, \emph{\textquotedblleft{}short.\textquotedblright{}}
\end{itemize}

You already own \emph{brevis} in English, several times over:

\begin{itemize}
\raggedright
  \item \textbf{brief} — short (a \emph{brief} meeting; a legal \emph{brief}).
  \item \textbf{brevity} — shortness (\textquotedblleft{}brevity is the soul of wit\textquotedblright{}).
  \item \textbf{abbreviate}, \textbf{abridge} — to make short.
  \item even the musical \textbf{breve} and the little \textbf{breve} mark ˘ over a short vowel.
\end{itemize}

So \emph{até breve} is \textbf{\textquotedblleft{}until {[}a{]} short {[}while{]}\textquotedblright{}} — see you before long. It's a touch warmer and less fixed than \emph{até logo}; both mean \textquotedblleft{}see you soon.\textquotedblright{}
\end{cousinweb}

\subsection*{Your turn}
Say these aloud:

\begin{itemize}
\raggedright
  \item \textquotedblleft{}até breve\textquotedblright{} — \emph{ah-TEH BREH-vee}
  \item \textquotedblleft{}breve\textquotedblright{} then English \textquotedblleft{}brief, brevity, abbreviate\textquotedblright{} — one root
  \item the three \emph{até} goodbyes — até logo / até amanhã / até breve
\end{itemize}

\begin{tcolorbox}[breakable,colback=teal!4,colframe=teal!35!black,title={Before you move on}]
What does \emph{breve} mean, and its English cousins? (\textquotedblleft{}Short\textquotedblright{} — brief, brevity, abbreviate.) So \emph{até breve} literally says? (\textquotedblleft{}Until a short while\textquotedblright{} $\to$ see you soon.) Next: put the Portuguese goodbyes together.
\end{tcolorbox}
\section[Practice]{Practice — saying goodbye in Portuguese}

\label{lesson:PT-C04-practice}



\begin{quote}
No new words. You can now greet, ask \emph{tudo bem?}, and \textbf{close} the conversation. Run the goodbyes until the right one is automatic.
\end{quote}

\subsection*{You'll want to know: The closings}
\textbf{Final / formal} — the \textquotedblleft{}to God\textquotedblright{} goodbye:

\begin{quote}
— Foi um prazer. \textbf{Adeus!}
\end{quote}

\textbf{Everyday \textquotedblleft{}see you\textquotedblright{}} — pick the \emph{até} that fits:

\begin{quote}
— \textbf{Até logo!} (see you later / soon) — \textbf{Até amanhã!} (see you tomorrow) — \textbf{Até breve!} (see you soon)
\end{quote}

\textbf{Casual chain} — greet and part:

\begin{quote}
— Olá! Tudo bem? … \textbf{Até logo!} — \textbf{Até!} (Brazilians often clip it to just \emph{\textquotedblleft{}até!\textquotedblright{}})
\end{quote}

\begin{grammarlens}[title={What you've built this chapter}]
\begin{itemize}
\raggedright
  \item \textbf{adeus} — goodbye (\textquotedblleft{}to God\textquotedblright{}; the \emph{adiós / adieu / addio} family, and English \emph{goodbye} = \textquotedblleft{}God be with ye\textquotedblright{}).
  \item \textbf{até logo} — see you later (\emph{até} $\leftarrow$ \textbf{Arabic \emph{ḥattā}}, twin of Spanish \emph{hasta}; \emph{logo} $\leftarrow$ \emph{locō}, kin of \emph{luego}). \emph{Até logo} = \emph{hasta luego}.
  \item \textbf{até amanhã} — see you tomorrow (\emph{amanhã} $\leftarrow$ \emph{ad māneāna}, kin of \emph{mañana}; the \emph{nh} = Spanish \emph{ñ}).
  \item \textbf{até breve} — see you soon (\emph{breve} $\leftarrow$ \emph{brevis}, \textquotedblleft{}short\textquotedblright{} — brief/brevity).
\end{itemize}
\end{grammarlens}

\subsection*{Your turn}
Say these aloud:

\begin{itemize}
\raggedright
  \item greet, ask, and close — \textquotedblleft{}Olá! Tudo bem? … Até amanhã!\textquotedblright{}
  \item choose the right goodbye — leaving for good / till tomorrow / soon
  \item \emph{até logo} and Spanish \emph{hasta luego} — twins, Arabic \emph{até/hasta} and all
\end{itemize}

\emph{Twice through:} Open and close a whole tiny exchange without stopping.

\begin{tcolorbox}[breakable,colback=teal!4,colframe=teal!35!black,title={Before you move on}]
Which Portuguese goodbye means \textquotedblleft{}to God\textquotedblright{}? (\emph{Adeus}.) Which word inside the \emph{até} goodbyes came from Arabic, and what's its Spanish twin? (\emph{Até} $\leftarrow$ \emph{ḥattā} — Spanish \emph{hasta}.) You can now run a Portuguese conversation start to finish — the full greet-to-goodbye arc.
\end{tcolorbox}
